\hypertarget{appendix-cc0909a7ef14}{}
\section*{Page faults and demand paging}
\section{Page faults and demand paging}\label{page-faults-and-demand-paging}

\textbf{Curriculum depth:} Introduced · \textbf{First seen:} C003 · \textbf{Latest development:} C003

\subsection{Definition and mechanism}\label{definition-and-mechanism}

A memory fault requires OS intervention when current translation/protection state cannot satisfy access. Valid demand loading, zero-fill, or copy-on-write may install/alter mapping then retry. Invalid addresses or denied permissions may fail instead. A fault need not perform disk I/O; a TLB miss need not fault at all.

\subsection{Prerequisites}\label{prerequisites}

\begin{itemize}
\tightlist
\item
  \hyperref[appendix-01e5905165c4]{Page tables and MMU} --- Cached or faulting translation starts from the mapping/protection mechanism.
\end{itemize}

\subsection{Used by}\label{used-by}

\begin{itemize}
\tightlist
\item
  \hyperref[appendix-7a06e5baac04]{Copy-on-write} --- uses this mechanism as stated in its prerequisites.
\item
  \hyperref[appendix-313ee7a0ba7b]{Memory-mapped files} --- uses this mechanism as stated in its prerequisites.
\end{itemize}

\subsection{Recall and next depth}\label{recall-and-next-depth}

Trace successful first-touch recovery versus denied access.

\textbf{Pending:} Minor/major faults, working sets, swapping, and platform terminology.

\subsection{Source and navigation}\label{source-and-navigation}

\href{Continuous\%20Notes/Computing\%20Foundations\%20-\%20From\%20Bits\%20to\%20C++\%20and\%20Systems\%20APIs\#Where\%20the\%20bytes\%20live}{C003 teaching chapter} · \hyperref[appendix-fd090d982a15]{Part definitions} · \href{}{Dependency Map (02-Foundations-II, appendix-4ad47bdd269a)}
